%% ============================================================
%%  Section: Data and Sample Construction
%%  Depositor Discipline in Modern Banking (2026-04-06 draft)
%%
%%  Scope: data provenance, variable construction, winsorization,
%%  depositor sophistication index, large-bank validation sample.
%%  Treatment variable formulas are in sec:background (eq:dayone,
%%  eq:treatment, eq:binary). Sample size and balance table are
%%  in the descriptive statistics section that follows.
%%
%%  NOTE: The desc-stats section currently carries \label{sec:data}.
%%  When incorporating this file, change the desc-stats section
%%  heading to \section{Descriptive Statistics}\label{sec:descriptive}
%%  to avoid duplicate labels.
%%
%%  New citation needed: narayanan2026decline (see bib note below).
%%  New citation needed: IRS_SOI (bib key placeholder used below).
%% ============================================================

\section{Data and Sample Construction}
\label{sec:data}

\subsection{Data Sources}
\label{sec:data:sources}

Our empirical design draws on four data sources. The primary source is quarterly
Call Reports (FFIEC 041 and 051) filed by all federally insured commercial banks
and savings institutions with the Federal Financial Institutions Examination
Council (FFIEC). Call Reports provide balance-sheet stocks and income-statement
flows, including deposit-account breakdowns by insurance status. Beginning with each
institution's CECL adoption quarter, the cumulative-effect retained-earnings
adjustment recorded in Schedule~RI-A. This last field is the main input for the
treatment variable (\texttt{CECL Adj}) defined in Section~\ref{sec:background}.

For the depositor sophistication analysis, we supplement the Call Report panel
with two additional cross-sections. The 2022 FDIC Summary of Deposits (SOD)
provides deposit balances at the branch level, reported as of June~30 each year,
for all FDIC-insured institutions. We merge these branch-level balances to
five-year ZIP-code estimates from the 2022 American Community Survey and IRS
Statistics of Income (SOI) ZIP-code tabulations to construct a
bank-level measure of depositor financial sophistication
(Section~\ref{sec:data:sophistication}).

The signal-validation exercise requires market data available only for large
publicly traded institutions. We use FR~Y-9C consolidated holding-company
financial statements filed with the Federal Reserve, monthly equity market
capitalization from CRSP, and five-year credit default swap (CDS) spreads from
Bloomberg. CRSP permanent company identifiers are linked to Federal Reserve RSSD
identifiers via the CRSP--FRB crosswalk maintained by the Federal Reserve Bank
of New York.

\subsection{Call Report Panel Construction}
\label{sec:data:panel}

We construct a bank-quarter panel identified by each institution's RSSD identifier
and quarter, spanning 2016Q1--2025Q4 for all institutions. Income-statement
variables reported on a year-to-date basis (interest expense, interest income,
net interest income, net income, noninterest expense, and noninterest income) are
converted to quarterly flows by differencing within bank-year.

CECL adoption is identified by searching the window 2022Q3--2023Q4 for the
first quarter in which a bank reports a non-missing value for the
cumulative-effect retained-earnings adjustment in Schedule~RI-A. This quarter is
assigned as the bank-specific adoption date $t^{*}$, and the treatment variable
\texttt{CECL Adj} is set to its adoption-quarter value and held fixed
throughout the panel (Section~\ref{sec:background}, equation~\eqref{eq:treatment}).
Banks with no non-missing adjustment in the eligible window are excluded. The
community-bank sample is further restricted to institutions with total assets
below \$10~billion as of 2022Q4, consistent with the Federal Reserve's standard
definition; the large SEC-filing institutions required to adopt in 2020Q1 are
excluded from the main panel and analyzed separately in the signal-validation
exercise.

\subsection{Outcome Variables}
\label{sec:data:outcomes}

We construct two groups of outcome variables, expressed as percentages of
contemporaneous total assets and measured at quarterly frequency.

The primary deposit outcomes are uninsured time deposits and insured time
deposits, each reported directly in Schedule~RC-E of the Call Report and scaled
by total assets. Time deposits are particularly well suited to detecting
depositor discipline since a decline in balances reflects a deliberate choice not 
to renew rather than routine liquidity management.

The funding-cost outcome is quarterly interest expense scaled by average total
assets; net interest margin is quarterly net interest income scaled by average
total assets. Profitability is measured by return on assets, computed as
quarterly net income to average total assets.

\subsection{Control Variables and Winsorization}
\label{sec:data:controls}

All regression specifications include two lagged balance-sheet controls measured
as of the quarter prior to the outcome: the natural logarithm of total assets
and the equity-to-assets ratio. Deposit specifications additionally include lagged interest expense scaled by
total assets; this control is excluded from interest expense, net interest
margin, and return on assets specifications to avoid mechanical correlation with
the dependent variable. Log assets accommodates scale effects on deposit composition. The
equity ratio controls for the bank's capital position. Lagged interest
expense captures pre-existing differences in deposit pricing strategy and, to the
extent that banks anticipated the CECL shock and adjusted deposit rates ahead of
formal adoption, absorbs any such anticipatory response. All ratio variables, both
outcomes and controls, are winsorized within each calendar quarter at the 1st and
99th percentiles.

\subsection{Depositor Sophistication}
\label{sec:data:sophistication}

We construct a bank-level measure of depositor financial sophistication by
merging branch-level deposit data with ZIP-code demographics. Following
\citet{narayanan2026decline}, we classify ZIP codes as financially sophisticated
using the 2022 American Community Survey and IRS SOI tabulations: a ZIP code is
assigned a sophistication indicator of one if its share of adults with at least a
bachelor's degree \textit{and} its share of tax filers reporting investment
income both exceed their respective sample medians, and zero otherwise.
For each bank~$i$, we aggregate to the bank level using deposit weights from the
2022 FDIC Summary of Deposits:
\begin{equation}
  \texttt{Sophisticated}_{i}
  \;=\;
  \frac{
    \displaystyle\sum_{b \in \mathcal{B}_i}
      \mathrm{Deposits}_{b} \cdot \mathbf{1}\!\left\{S_{b} = 1\right\}
  }{
    \displaystyle\sum_{b \in \mathcal{B}_i} \mathrm{Deposits}_{b}
  },
  \label{eq:sophistication}
\end{equation}
where $\mathcal{B}_i$ is the set of branches belonging to bank~$i$,
$\mathrm{Deposits}_b$ is branch~$b$'s June~2022 deposit balance, and $S_b$ is
the ZIP-code sophistication indicator for branch~$b$'s location.
\texttt{Sophisticated}$_i$ is thus the deposit-weighted fraction of a bank's branch-level
balances located in high-education, high-investment-income ZIP codes. Banks are
classified as \emph{high sophistication} if their score is at or above the 75th
percentile of the cross-sectional distribution.

\subsection{Large-Bank Validation Sample}
\label{sec:data:validation}

To establish the information content of the CECL Day-One disclosure, we
construct a complementary panel of large, publicly traded bank holding companies
that adopted CECL in 2020Q1 under SEC-filer requirements
(Section~\ref{sec:identification:signal}). From FR~Y-9C filings, we identify
each institution's first CECL reporting quarter and compute the treatment
variable following equation~\eqref{eq:treatment}.

For the equity market analysis, we extract monthly market capitalization from
CRSP for each holding company, linked to FR~Y-9C via the CRSP--FRB crosswalk.
For the CDS analysis, we obtain monthly five-year CDS spreads from Bloomberg,
retaining only contracts with observable trading activity.

Although the validation sample differs from the community-bank panel in size,
listing status, and adoption timing, it shares the same treatment concept,
cross-sectional variation in the Day-One CECL adjustment relative to equity,
and the same event-time structure. Its purpose is to verify that the CECL
disclosure is an information event rather than a mechanical accounting
reclassification, before testing depositor responses in the community-bank
setting.
